% Options for packages loaded elsewhere
\PassOptionsToPackage{unicode}{hyperref}
\PassOptionsToPackage{hyphens}{url}
\PassOptionsToPackage{dvipsnames,svgnames,x11names}{xcolor}
%
\documentclass[
  letterpaper,
  DIV=11,
  numbers=noendperiod]{scrartcl}

\usepackage{amsmath,amssymb}
\usepackage{iftex}
\ifPDFTeX
  \usepackage[T1]{fontenc}
  \usepackage[utf8]{inputenc}
  \usepackage{textcomp} % provide euro and other symbols
\else % if luatex or xetex
  \usepackage{unicode-math}
  \defaultfontfeatures{Scale=MatchLowercase}
  \defaultfontfeatures[\rmfamily]{Ligatures=TeX,Scale=1}
\fi
\usepackage{lmodern}
\ifPDFTeX\else  
    % xetex/luatex font selection
\fi
% Use upquote if available, for straight quotes in verbatim environments
\IfFileExists{upquote.sty}{\usepackage{upquote}}{}
\IfFileExists{microtype.sty}{% use microtype if available
  \usepackage[]{microtype}
  \UseMicrotypeSet[protrusion]{basicmath} % disable protrusion for tt fonts
}{}
\makeatletter
\@ifundefined{KOMAClassName}{% if non-KOMA class
  \IfFileExists{parskip.sty}{%
    \usepackage{parskip}
  }{% else
    \setlength{\parindent}{0pt}
    \setlength{\parskip}{6pt plus 2pt minus 1pt}}
}{% if KOMA class
  \KOMAoptions{parskip=half}}
\makeatother
\usepackage{xcolor}
\setlength{\emergencystretch}{3em} % prevent overfull lines
\setcounter{secnumdepth}{-\maxdimen} % remove section numbering
% Make \paragraph and \subparagraph free-standing
\makeatletter
\ifx\paragraph\undefined\else
  \let\oldparagraph\paragraph
  \renewcommand{\paragraph}{
    \@ifstar
      \xxxParagraphStar
      \xxxParagraphNoStar
  }
  \newcommand{\xxxParagraphStar}[1]{\oldparagraph*{#1}\mbox{}}
  \newcommand{\xxxParagraphNoStar}[1]{\oldparagraph{#1}\mbox{}}
\fi
\ifx\subparagraph\undefined\else
  \let\oldsubparagraph\subparagraph
  \renewcommand{\subparagraph}{
    \@ifstar
      \xxxSubParagraphStar
      \xxxSubParagraphNoStar
  }
  \newcommand{\xxxSubParagraphStar}[1]{\oldsubparagraph*{#1}\mbox{}}
  \newcommand{\xxxSubParagraphNoStar}[1]{\oldsubparagraph{#1}\mbox{}}
\fi
\makeatother

\usepackage{color}
\usepackage{fancyvrb}
\newcommand{\VerbBar}{|}
\newcommand{\VERB}{\Verb[commandchars=\\\{\}]}
\DefineVerbatimEnvironment{Highlighting}{Verbatim}{commandchars=\\\{\}}
% Add ',fontsize=\small' for more characters per line
\usepackage{framed}
\definecolor{shadecolor}{RGB}{241,243,245}
\newenvironment{Shaded}{\begin{snugshade}}{\end{snugshade}}
\newcommand{\AlertTok}[1]{\textcolor[rgb]{0.68,0.00,0.00}{#1}}
\newcommand{\AnnotationTok}[1]{\textcolor[rgb]{0.37,0.37,0.37}{#1}}
\newcommand{\AttributeTok}[1]{\textcolor[rgb]{0.40,0.45,0.13}{#1}}
\newcommand{\BaseNTok}[1]{\textcolor[rgb]{0.68,0.00,0.00}{#1}}
\newcommand{\BuiltInTok}[1]{\textcolor[rgb]{0.00,0.23,0.31}{#1}}
\newcommand{\CharTok}[1]{\textcolor[rgb]{0.13,0.47,0.30}{#1}}
\newcommand{\CommentTok}[1]{\textcolor[rgb]{0.37,0.37,0.37}{#1}}
\newcommand{\CommentVarTok}[1]{\textcolor[rgb]{0.37,0.37,0.37}{\textit{#1}}}
\newcommand{\ConstantTok}[1]{\textcolor[rgb]{0.56,0.35,0.01}{#1}}
\newcommand{\ControlFlowTok}[1]{\textcolor[rgb]{0.00,0.23,0.31}{\textbf{#1}}}
\newcommand{\DataTypeTok}[1]{\textcolor[rgb]{0.68,0.00,0.00}{#1}}
\newcommand{\DecValTok}[1]{\textcolor[rgb]{0.68,0.00,0.00}{#1}}
\newcommand{\DocumentationTok}[1]{\textcolor[rgb]{0.37,0.37,0.37}{\textit{#1}}}
\newcommand{\ErrorTok}[1]{\textcolor[rgb]{0.68,0.00,0.00}{#1}}
\newcommand{\ExtensionTok}[1]{\textcolor[rgb]{0.00,0.23,0.31}{#1}}
\newcommand{\FloatTok}[1]{\textcolor[rgb]{0.68,0.00,0.00}{#1}}
\newcommand{\FunctionTok}[1]{\textcolor[rgb]{0.28,0.35,0.67}{#1}}
\newcommand{\ImportTok}[1]{\textcolor[rgb]{0.00,0.46,0.62}{#1}}
\newcommand{\InformationTok}[1]{\textcolor[rgb]{0.37,0.37,0.37}{#1}}
\newcommand{\KeywordTok}[1]{\textcolor[rgb]{0.00,0.23,0.31}{\textbf{#1}}}
\newcommand{\NormalTok}[1]{\textcolor[rgb]{0.00,0.23,0.31}{#1}}
\newcommand{\OperatorTok}[1]{\textcolor[rgb]{0.37,0.37,0.37}{#1}}
\newcommand{\OtherTok}[1]{\textcolor[rgb]{0.00,0.23,0.31}{#1}}
\newcommand{\PreprocessorTok}[1]{\textcolor[rgb]{0.68,0.00,0.00}{#1}}
\newcommand{\RegionMarkerTok}[1]{\textcolor[rgb]{0.00,0.23,0.31}{#1}}
\newcommand{\SpecialCharTok}[1]{\textcolor[rgb]{0.37,0.37,0.37}{#1}}
\newcommand{\SpecialStringTok}[1]{\textcolor[rgb]{0.13,0.47,0.30}{#1}}
\newcommand{\StringTok}[1]{\textcolor[rgb]{0.13,0.47,0.30}{#1}}
\newcommand{\VariableTok}[1]{\textcolor[rgb]{0.07,0.07,0.07}{#1}}
\newcommand{\VerbatimStringTok}[1]{\textcolor[rgb]{0.13,0.47,0.30}{#1}}
\newcommand{\WarningTok}[1]{\textcolor[rgb]{0.37,0.37,0.37}{\textit{#1}}}

\providecommand{\tightlist}{%
  \setlength{\itemsep}{0pt}\setlength{\parskip}{0pt}}\usepackage{longtable,booktabs,array}
\usepackage{calc} % for calculating minipage widths
% Correct order of tables after \paragraph or \subparagraph
\usepackage{etoolbox}
\makeatletter
\patchcmd\longtable{\par}{\if@noskipsec\mbox{}\fi\par}{}{}
\makeatother
% Allow footnotes in longtable head/foot
\IfFileExists{footnotehyper.sty}{\usepackage{footnotehyper}}{\usepackage{footnote}}
\makesavenoteenv{longtable}
\usepackage{graphicx}
\makeatletter
\def\maxwidth{\ifdim\Gin@nat@width>\linewidth\linewidth\else\Gin@nat@width\fi}
\def\maxheight{\ifdim\Gin@nat@height>\textheight\textheight\else\Gin@nat@height\fi}
\makeatother
% Scale images if necessary, so that they will not overflow the page
% margins by default, and it is still possible to overwrite the defaults
% using explicit options in \includegraphics[width, height, ...]{}
\setkeys{Gin}{width=\maxwidth,height=\maxheight,keepaspectratio}
% Set default figure placement to htbp
\makeatletter
\def\fps@figure{htbp}
\makeatother

\KOMAoption{captions}{tableheading}
\makeatletter
\@ifpackageloaded{caption}{}{\usepackage{caption}}
\AtBeginDocument{%
\ifdefined\contentsname
  \renewcommand*\contentsname{Table of contents}
\else
  \newcommand\contentsname{Table of contents}
\fi
\ifdefined\listfigurename
  \renewcommand*\listfigurename{List of Figures}
\else
  \newcommand\listfigurename{List of Figures}
\fi
\ifdefined\listtablename
  \renewcommand*\listtablename{List of Tables}
\else
  \newcommand\listtablename{List of Tables}
\fi
\ifdefined\figurename
  \renewcommand*\figurename{Figure}
\else
  \newcommand\figurename{Figure}
\fi
\ifdefined\tablename
  \renewcommand*\tablename{Table}
\else
  \newcommand\tablename{Table}
\fi
}
\@ifpackageloaded{float}{}{\usepackage{float}}
\floatstyle{ruled}
\@ifundefined{c@chapter}{\newfloat{codelisting}{h}{lop}}{\newfloat{codelisting}{h}{lop}[chapter]}
\floatname{codelisting}{Listing}
\newcommand*\listoflistings{\listof{codelisting}{List of Listings}}
\makeatother
\makeatletter
\makeatother
\makeatletter
\@ifpackageloaded{caption}{}{\usepackage{caption}}
\@ifpackageloaded{subcaption}{}{\usepackage{subcaption}}
\makeatother

\ifLuaTeX
  \usepackage{selnolig}  % disable illegal ligatures
\fi
\usepackage{bookmark}

\IfFileExists{xurl.sty}{\usepackage{xurl}}{} % add URL line breaks if available
\urlstyle{same} % disable monospaced font for URLs
\hypersetup{
  pdftitle={Simulation of Crop production function using AquaCrop-OSPy},
  colorlinks=true,
  linkcolor={blue},
  filecolor={Maroon},
  citecolor={Blue},
  urlcolor={Blue},
  pdfcreator={LaTeX via pandoc}}


\title{Simulation of Crop production function using AquaCrop-OSPy}
\author{}
\date{}

\begin{document}
\maketitle


This README file provides users with an overview of how we use
AquaCrop-OSPy---an open-source Python implementation of the U.N. Food
and Agriculture Organization (FAO) AquaCrop model---to simulate crop
production functions in the Lake Diefenbaker area

Please note that these codes are designed to run in a Python
environment. Ensure that you have a Python setup in place and that all
required packages are installed before importing them.

\begin{verbatim}
Using virtual environment '/Users/tharakajayalath/.virtualenvs/r-reticulate' ...
\end{verbatim}

\begin{Shaded}
\begin{Highlighting}[]
\CommentTok{\# !pip install git+https://github.com/aquacropos/aquacrop}

\CommentTok{\# from google.colab import output}
\CommentTok{\# output.clear()}
\end{Highlighting}
\end{Shaded}

Now that AquaCrop is installed, we need to import its components into
the environment.

\begin{Shaded}
\begin{Highlighting}[]
\ImportTok{from}\NormalTok{ aquacrop }\ImportTok{import}\NormalTok{ AquaCropModel, Soil, Crop, InitialWaterContent, FieldMngt, GroundWater}
\ImportTok{from}\NormalTok{ aquacrop.utils }\ImportTok{import}\NormalTok{ prepare\_weather, get\_filepath}

\ImportTok{import}\NormalTok{ pandas }\ImportTok{as}\NormalTok{ pd}
\ImportTok{import}\NormalTok{ matplotlib.pyplot }\ImportTok{as}\NormalTok{ plt}
\ImportTok{import}\NormalTok{ seaborn }\ImportTok{as}\NormalTok{ sns}
\end{Highlighting}
\end{Shaded}

Running an AquaCrop-OSPy model requires selecting five key components:

\begin{enumerate}
\def\labelenumi{\arabic{enumi}.}
\tightlist
\item
  Daily climate measurements
\item
  Soil selection
\item
  Crop selection
\item
  Initial water content
\item
  Simulation start and end dates
\end{enumerate}

\textbf{1. Climate measurements}

AquaCrop-OSPy requires weather data for the cropping period being
simulated. This includes a daily time series of minimum and maximum
temperatures (°C), precipitation (mm), and reference crop
evapotranspiration (ET₀, mm).

\textbf{2. Soil}

Selecting a soil type for an AquaCrop-OSPy simulation is done via the
Soil object. This object contains all the compositional and hydraulic
properties required for the model. The simplest way to choose a soil is
to use one of the built-in soil types provided in AquaCrop defaults.
Users can modify any of these built-in soil types or create their own
custom soil profiles.

\textbf{3. Crop}

The crop type used in the simulation is selected in a similar way via
the Crop object. To define a crop, you must specify the crop type and
planting date. Any of the built-in crop types (currently Maize, Wheat,
Rice, and Potato) can be selected or user can define crop parameters.
Users may also edit the built-in crop parameters or create custom crop
definitions.

Broadly these crop parameters are divide into two categories:

a). Conservative parameters

b). Non-conservative parameters - Needs to calibrate based on local
conditions

\textbf{4. Initial water content}

The initial soil-water content at the beginning of the simulation is
specified using the \emph{InitialWaterContent} object.

\textbf{5. Simulation start and end dates}

You must define the simulation's start and end dates (YYYY/MM/DD). These
determine the period over which AquaCrop-OSPy will run.

\textbf{Running the simulation}

Additional model components may also be specified, including irrigation
management, field management, groundwater conditions. Once you have
defined the weather data, crop, soil, and initial water content, you are
ready to run the simulation. To run the model, combine all components
into an AquaCropModel object. This is also where you specify the
simulation start and end dates.

\textbf{AquaCrop-OSPy Simulation Configuration to Simulate Crop
Yield--Water Relationship for Lake Diefenbaker}

Now that you have a basic understanding of how AquaCrop-OSPy works, I
will walk you through a real example in which we simulate the crop
production function for wheat using the actual data employed in our
manuscript. In our simulation setup for the Lake Diefenbaker region, as
used in the manuscript, there are a few details that need to be
clarified.

a). In our case we simulate 394 sampling sites not only 1 site

b). We simulate crop production during 2018 - 2023

First, I will explain the simulation settings for revealing the crop
yield--water relationship, which we later use to estimate the average
value of irrigation water. Here, we calculate the average value as:

\begin{equation}

{AV_{itj}} =     \sum_{j=1}^{n} \frac{(Y^{IR}_{itj} -Y^{RF}_{itj})*P_{it}-(C^{IR}_{it}-C^{RF}_{it})}{IRWQ_{itj}}
\label{equation1}
\end{equation}

Where;

\(AV_{ijt}=\) Average value of unit of irrigation water used in
\(j^{th}\) field for \(i^{th}\) crop during the \(t^{th}\) year

\(Y^{IR}_{itj}=\) Yield from \(i^{th}\) crop in year \(t\) at the
\(j^{th}\) field with net irrigation demand

\(Y^{RF}_{itj}=\) Yield from \(i^{th}\) crop in year \(t\) at the
\(j^{th}\) field under rain-fed

\(P_{it}\) = \(i^{th}\) Crop price during \(t^{th}\) year

\(C^{IR}_{it},C^{RF}_{it}=\) Operational cost of \(i^{th}\) crop under
irrigated and rain-fed production system during year \(t^{th}\)
respectively

\(IRWQ_{itj}=\) Total amount of irrigation water used in \(i^{th}\) crop
at \(j^{th}\) site during the \(t^{th}\) year

Here, the aim of the AquaCrop simulation is to model crop yields for
wheat under rain-fed conditions, as well as under irrigated conditions,
and to estimate the net irrigation requirement needed to achieve their
maximum potential yield.

\textbf{Simulation Configuration 1}

Here, I am referring to our first simulation configuration, where we
simulate the crop production function under rainfed conditions. The
following illustrative code is for the simulation of wheat production,
and later I will discuss which parameters need to be changed for
different crops.

\textbf{Step 1:}

Developing the weather data file in our case requires an additional step
because the Daymet and CMIP weather data files do not include ET₀.
Therefore, we need to calculate ET₀ based on the available weather
variables from the original file and then create the weather data file
accordingly. Each step in the code includes a note to help you
understand the process easily.

\begin{Shaded}
\begin{Highlighting}[]
\ImportTok{import}\NormalTok{ os}
\ImportTok{import}\NormalTok{ pandas }\ImportTok{as}\NormalTok{ pd}
\ImportTok{import}\NormalTok{ numpy }\ImportTok{as}\NormalTok{ np}

\CommentTok{\# {-}{-}{-} Paths {-}{-}{-}}
\CommentTok{\# Working directory}
\NormalTok{cwd }\OperatorTok{=}\NormalTok{ os.getcwd()}
\BuiltInTok{print}\NormalTok{(}\SpecialStringTok{f"Current working directory: }\SpecialCharTok{\{}\NormalTok{cwd}\SpecialCharTok{\}}\SpecialStringTok{"}\NormalTok{)}
\end{Highlighting}
\end{Shaded}

\begin{verbatim}
Current working directory: /Users/tharakajayalath/Library/CloudStorage/OneDrive-UniversityofSaskatchewan/Chapter II-IrrigationValue/Chapter-II/Codes
\end{verbatim}

\begin{Shaded}
\begin{Highlighting}[]
\CommentTok{\# Output directory}
\NormalTok{output\_dir }\OperatorTok{=}\NormalTok{ os.path.join(cwd, }\StringTok{\textquotesingle{}AquaCropReadme\textquotesingle{}}\NormalTok{)}
\NormalTok{os.makedirs(output\_dir, exist\_ok}\OperatorTok{=}\VariableTok{True}\NormalTok{)  }\CommentTok{\# Create directory if it doesn\textquotesingle{}t exist}

\CommentTok{\# Input file path}
\NormalTok{input\_file }\OperatorTok{=} \StringTok{\textquotesingle{}/Users/tharakajayalath/Library/CloudStorage/OneDrive{-}UniversityofSaskatchewan/Chapter II{-}IrrigationValue/Chapter{-}II/AquaCropOPSyData/ClimateData/weather\_data\_Aqua.csv\textquotesingle{}}

\CommentTok{\# Output file path}
\NormalTok{output\_file }\OperatorTok{=}\NormalTok{ os.path.join(output\_dir, }\StringTok{\textquotesingle{}daymet\_data\_with\_et0.csv\textquotesingle{}}\NormalTok{)}

\CommentTok{\# {-}{-}{-} Load Daymet data {-}{-}{-}}
\NormalTok{climate\_data }\OperatorTok{=}\NormalTok{ pd.read\_csv(input\_file, on\_bad\_lines}\OperatorTok{=}\StringTok{\textquotesingle{}skip\textquotesingle{}}\NormalTok{)}

\CommentTok{\# {-}{-}{-} Check Date column and process {-}{-}{-}}
\ControlFlowTok{if} \StringTok{\textquotesingle{}Date\textquotesingle{}} \KeywordTok{in}\NormalTok{ climate\_data.columns:}
\NormalTok{    climate\_data[}\StringTok{\textquotesingle{}Date\textquotesingle{}}\NormalTok{] }\OperatorTok{=}\NormalTok{ pd.to\_datetime(climate\_data[}\StringTok{\textquotesingle{}Date\textquotesingle{}}\NormalTok{], errors}\OperatorTok{=}\StringTok{\textquotesingle{}coerce\textquotesingle{}}\NormalTok{)}
\NormalTok{    climate\_data }\OperatorTok{=}\NormalTok{ climate\_data.dropna(subset}\OperatorTok{=}\NormalTok{[}\StringTok{\textquotesingle{}Date\textquotesingle{}}\NormalTok{])}
    
    \CommentTok{\# Extract Day, Month, Year}
\NormalTok{    climate\_data[}\StringTok{\textquotesingle{}Day\textquotesingle{}}\NormalTok{] }\OperatorTok{=}\NormalTok{ climate\_data[}\StringTok{\textquotesingle{}Date\textquotesingle{}}\NormalTok{].dt.day}
\NormalTok{    climate\_data[}\StringTok{\textquotesingle{}Month\textquotesingle{}}\NormalTok{] }\OperatorTok{=}\NormalTok{ climate\_data[}\StringTok{\textquotesingle{}Date\textquotesingle{}}\NormalTok{].dt.month}
\NormalTok{    climate\_data[}\StringTok{\textquotesingle{}Year\textquotesingle{}}\NormalTok{] }\OperatorTok{=}\NormalTok{ climate\_data[}\StringTok{\textquotesingle{}Date\textquotesingle{}}\NormalTok{].dt.year}

    \CommentTok{\# {-}{-}{-} Calculate ET0 {-}{-}{-}}
\NormalTok{    T\_max }\OperatorTok{=}\NormalTok{ climate\_data[}\StringTok{\textquotesingle{}MaxTemp\textquotesingle{}}\NormalTok{]}
\NormalTok{    T\_min }\OperatorTok{=}\NormalTok{ climate\_data[}\StringTok{\textquotesingle{}MinTemp\textquotesingle{}}\NormalTok{]}
\NormalTok{    precipitation }\OperatorTok{=}\NormalTok{ climate\_data[}\StringTok{\textquotesingle{}Precipitation\textquotesingle{}}\NormalTok{]}
\NormalTok{    solar\_radiation }\OperatorTok{=}\NormalTok{ climate\_data[}\StringTok{\textquotesingle{}R\_n\textquotesingle{}}\NormalTok{]}
\NormalTok{    e\_a }\OperatorTok{=}\NormalTok{ climate\_data[}\StringTok{\textquotesingle{}e\_a\textquotesingle{}}\NormalTok{]}

\NormalTok{    T\_mean }\OperatorTok{=}\NormalTok{ (T\_max }\OperatorTok{+}\NormalTok{ T\_min) }\OperatorTok{/} \DecValTok{2}
\NormalTok{    e\_s }\OperatorTok{=} \FloatTok{0.6108} \OperatorTok{*}\NormalTok{ np.exp((}\FloatTok{17.27} \OperatorTok{*}\NormalTok{ T\_mean) }\OperatorTok{/}\NormalTok{ (T\_mean }\OperatorTok{+} \FloatTok{237.3}\NormalTok{))}
\NormalTok{    R\_n }\OperatorTok{=}\NormalTok{ solar\_radiation}
\NormalTok{    G }\OperatorTok{=} \DecValTok{0}
\NormalTok{    gamma }\OperatorTok{=} \FloatTok{0.066}
\NormalTok{    u }\OperatorTok{=} \FloatTok{1.2}
\NormalTok{    delta }\OperatorTok{=}\NormalTok{ (}\DecValTok{4098} \OperatorTok{*}\NormalTok{ e\_s) }\OperatorTok{/}\NormalTok{ ((T\_mean }\OperatorTok{+} \FloatTok{237.3}\NormalTok{) }\OperatorTok{**} \DecValTok{2}\NormalTok{)}

\NormalTok{    ReferenceET }\OperatorTok{=}\NormalTok{ (}\FloatTok{0.408} \OperatorTok{*}\NormalTok{ delta }\OperatorTok{*}\NormalTok{ (R\_n }\OperatorTok{{-}}\NormalTok{ G) }\OperatorTok{+}\NormalTok{ gamma }\OperatorTok{*}\NormalTok{ (}\DecValTok{900} \OperatorTok{/}\NormalTok{ (T\_mean }\OperatorTok{+} \DecValTok{273}\NormalTok{)) }\OperatorTok{*}\NormalTok{ u }\OperatorTok{*}\NormalTok{ (e\_s }\OperatorTok{{-}}\NormalTok{ e\_a)) }\OperatorTok{/}\NormalTok{ (delta }\OperatorTok{+}\NormalTok{ gamma }\OperatorTok{*}\NormalTok{ (}\DecValTok{1} \OperatorTok{+} \FloatTok{0.34} \OperatorTok{*}\NormalTok{ u))}
\NormalTok{    climate\_data[}\StringTok{\textquotesingle{}ReferenceET\textquotesingle{}}\NormalTok{] }\OperatorTok{=}\NormalTok{ ReferenceET}

    \CommentTok{\# Reorder columns}
\NormalTok{    cols }\OperatorTok{=}\NormalTok{ [}\StringTok{\textquotesingle{}Day\textquotesingle{}}\NormalTok{, }\StringTok{\textquotesingle{}Month\textquotesingle{}}\NormalTok{, }\StringTok{\textquotesingle{}Year\textquotesingle{}}\NormalTok{] }\OperatorTok{+}\NormalTok{ [col }\ControlFlowTok{for}\NormalTok{ col }\KeywordTok{in}\NormalTok{ climate\_data.columns }\ControlFlowTok{if}\NormalTok{ col }\KeywordTok{not} \KeywordTok{in}\NormalTok{ [}\StringTok{\textquotesingle{}Day\textquotesingle{}}\NormalTok{, }\StringTok{\textquotesingle{}Month\textquotesingle{}}\NormalTok{, }\StringTok{\textquotesingle{}Year\textquotesingle{}}\NormalTok{]]}
\NormalTok{    climate\_data }\OperatorTok{=}\NormalTok{ climate\_data[cols]}

    \CommentTok{\# {-}{-}{-} Save CSV safely {-}{-}{-}}
\NormalTok{    climate\_data.to\_csv(output\_file, index}\OperatorTok{=}\VariableTok{False}\NormalTok{)}
    \BuiltInTok{print}\NormalTok{(}\SpecialStringTok{f"Reference ET₀ calculated and saved successfully to: }\SpecialCharTok{\{}\NormalTok{output\_file}\SpecialCharTok{\}}\SpecialStringTok{"}\NormalTok{)}
\ControlFlowTok{else}\NormalTok{:}
    \BuiltInTok{print}\NormalTok{(}\StringTok{"Error: \textquotesingle{}Date\textquotesingle{} column does not exist in the dataset."}\NormalTok{)}
\end{Highlighting}
\end{Shaded}

\begin{verbatim}
Reference ET₀ calculated and saved successfully to: /Users/tharakajayalath/Library/CloudStorage/OneDrive-UniversityofSaskatchewan/Chapter II-IrrigationValue/Chapter-II/Codes/AquaCropReadme/daymet_data_with_et0.csv
\end{verbatim}

\textbf{Step 2}

Run the AquaCrop-OSPy model for wheat across 394 sample sites within the
study area for the year 2018.

A few points need to be highlighted: for wheat, we use the
AquaCrop-defined conservative parameters and rely on published papers
from semi-arid regions to update the non-conservative parameters.

Please refer to the following code chunk to see how crop parameters are
handled using the Crop object. A complete list of crop parameters and
their descriptions is provided in the Appendix.

Also, note that we have controlled the irr\_mngt object. For rain-fed
conditions (no irrigation), we set it as:

irr\_mngt = IrrigationManagement(irrigation\_method=0)

For simulations under irrigation, we set it as:

irr\_mngt = IrrigationManagement(irrigation\_method=4, NetIrrSMT=70)

There are a few other potential irrigation setups, and the full list is
included in the Appendix.

Additional steps include looping the code across all 394 sites,
recording results for each site separately, and finally aggregating the
subset of simulated files for each site.

Here is the full code, and I have added some notes within it to explain
the aim of each code chunk.

\begin{Shaded}
\begin{Highlighting}[]
\CommentTok{\#\#\#2018{-}RAIN{-}FED\#\#018{-}RAIN{-}FED\#\#\#018{-}RAIN{-}FED\#\#\#018{-}RAIN{-}FED\#\#\#018{-}RAIN{-}FED\#\#\#018{-}RAIN{-}FED}

\ImportTok{import}\NormalTok{ os}
\ImportTok{import}\NormalTok{ pandas }\ImportTok{as}\NormalTok{ pd}
\ImportTok{import}\NormalTok{ numpy }\ImportTok{as}\NormalTok{ np}
\ImportTok{from}\NormalTok{ aquacrop }\ImportTok{import}\NormalTok{ AquaCropModel, Soil, Crop, InitialWaterContent, IrrigationManagement}
\ImportTok{from}\NormalTok{ aquacrop.utils }\ImportTok{import}\NormalTok{ prepare\_weather}

\CommentTok{\# {-}{-}{-} Paths {-}{-}{-}}
\NormalTok{cwd }\OperatorTok{=}\NormalTok{ os.getcwd()}
\BuiltInTok{print}\NormalTok{(}\SpecialStringTok{f"Current working directory: }\SpecialCharTok{\{}\NormalTok{cwd}\SpecialCharTok{\}}\SpecialStringTok{"}\NormalTok{)}
\end{Highlighting}
\end{Shaded}

\begin{Shaded}
\begin{Highlighting}[]
\CommentTok{\# Input weather data with ET0}
\NormalTok{input\_file }\OperatorTok{=} \StringTok{\textquotesingle{}./AquaCropReadme/daymet\_data\_with\_et0.csv\textquotesingle{}}

\CommentTok{\# Directories for weather files and simulation outputs}
\NormalTok{weather\_dir }\OperatorTok{=}\NormalTok{ os.path.join(cwd, }\StringTok{"ClimateData"}\NormalTok{)}
\NormalTok{sim\_output\_dir }\OperatorTok{=}\NormalTok{ os.path.join(cwd, }\StringTok{"SimulationResults"}\NormalTok{)}
\NormalTok{os.makedirs(weather\_dir, exist\_ok}\OperatorTok{=}\VariableTok{True}\NormalTok{)}\OperatorTok{;}
\NormalTok{os.makedirs(sim\_output\_dir, exist\_ok}\OperatorTok{=}\VariableTok{True}\NormalTok{)}\OperatorTok{;}


\CommentTok{\# Directory to save merged results}
\NormalTok{merged\_dir }\OperatorTok{=}\NormalTok{ os.path.join(cwd, }\StringTok{"AquaCropOPSyData/WheatRainfed"}\NormalTok{)}
\NormalTok{os.makedirs(merged\_dir, exist\_ok}\OperatorTok{=}\VariableTok{True}\NormalTok{)}\OperatorTok{;}
\NormalTok{merged\_output\_file }\OperatorTok{=}\NormalTok{ os.path.join(merged\_dir, }\StringTok{"wheat\_rainfed\_2018.csv"}\NormalTok{)}

\CommentTok{\# {-}{-}{-} Load data {-}{-}{-}}
\NormalTok{df }\OperatorTok{=}\NormalTok{ pd.read\_csv(input\_file)}
\NormalTok{unique\_sites }\OperatorTok{=}\NormalTok{ df[}\StringTok{\textquotesingle{}site\textquotesingle{}}\NormalTok{].unique()}
\BuiltInTok{print}\NormalTok{(}\SpecialStringTok{f"Total unique sites: }\SpecialCharTok{\{}\BuiltInTok{len}\NormalTok{(unique\_sites)}\SpecialCharTok{\}}\SpecialStringTok{"}\NormalTok{)}
\end{Highlighting}
\end{Shaded}

\begin{Shaded}
\begin{Highlighting}[]
\CommentTok{\# {-}{-}{-} Simulation dates {-}{-}{-}}
\NormalTok{sim\_start }\OperatorTok{=} \StringTok{\textquotesingle{}2018/01/01\textquotesingle{}}
\NormalTok{sim\_end }\OperatorTok{=} \StringTok{\textquotesingle{}2018/12/30\textquotesingle{}}

\CommentTok{\# {-}{-}{-} Loop through sites {-}{-}{-}}
\NormalTok{all\_simulation\_results }\OperatorTok{=}\NormalTok{ []}

\ControlFlowTok{for}\NormalTok{ site }\KeywordTok{in}\NormalTok{ unique\_sites:}
    \CommentTok{\# Filter data for site and year 2018}
\NormalTok{    df\_filtered }\OperatorTok{=}\NormalTok{ (df[df[}\StringTok{\textquotesingle{}site\textquotesingle{}}\NormalTok{] }\OperatorTok{==}\NormalTok{ site]}
\NormalTok{                   .loc[:, [}\StringTok{\textquotesingle{}Day\textquotesingle{}}\NormalTok{,}\StringTok{\textquotesingle{}Month\textquotesingle{}}\NormalTok{,}\StringTok{\textquotesingle{}Year\textquotesingle{}}\NormalTok{,}\StringTok{\textquotesingle{}MinTemp\textquotesingle{}}\NormalTok{,}\StringTok{\textquotesingle{}MaxTemp\textquotesingle{}}\NormalTok{,}\StringTok{\textquotesingle{}Precipitation\textquotesingle{}}\NormalTok{,}\StringTok{\textquotesingle{}ReferenceET\textquotesingle{}}\NormalTok{]]}
\NormalTok{                   .rename(columns}\OperatorTok{=}\NormalTok{\{}
                       \StringTok{\textquotesingle{}MinTemp\textquotesingle{}}\NormalTok{: }\StringTok{\textquotesingle{}Tmin(c)\textquotesingle{}}\NormalTok{,}
                       \StringTok{\textquotesingle{}MaxTemp\textquotesingle{}}\NormalTok{: }\StringTok{\textquotesingle{}Tmax(c)\textquotesingle{}}\NormalTok{,}
                       \StringTok{\textquotesingle{}Precipitation\textquotesingle{}}\NormalTok{: }\StringTok{\textquotesingle{}Prcp(mm)\textquotesingle{}}\NormalTok{,}
                       \StringTok{\textquotesingle{}ReferenceET\textquotesingle{}}\NormalTok{: }\StringTok{\textquotesingle{}Et0(mm)\textquotesingle{}}
\NormalTok{                   \})}
\NormalTok{                   .query(}\StringTok{"Year == 2018"}\NormalTok{))}
    
    \ControlFlowTok{if}\NormalTok{ df\_filtered.empty:}
        \BuiltInTok{print}\NormalTok{(}\SpecialStringTok{f"No data for site }\SpecialCharTok{\{}\NormalTok{site}\SpecialCharTok{\}}\SpecialStringTok{ in 2018. Skipping."}\NormalTok{)}
        \ControlFlowTok{continue}

    \CommentTok{\# Calculate total precipitation during growing period (May 1 {-} Oct 30)}
\NormalTok{    total\_precip }\OperatorTok{=}\NormalTok{ df\_filtered[(df\_filtered[}\StringTok{\textquotesingle{}Month\textquotesingle{}}\NormalTok{] }\OperatorTok{\textgreater{}=} \DecValTok{5}\NormalTok{) }\OperatorTok{\&}\NormalTok{ (df\_filtered[}\StringTok{\textquotesingle{}Month\textquotesingle{}}\NormalTok{] }\OperatorTok{\textless{}=} \DecValTok{10}\NormalTok{) }\OperatorTok{\&}
\NormalTok{                               ((df\_filtered[}\StringTok{\textquotesingle{}Month\textquotesingle{}}\NormalTok{] }\OperatorTok{!=} \DecValTok{10}\NormalTok{) }\OperatorTok{|}\NormalTok{ (df\_filtered[}\StringTok{\textquotesingle{}Day\textquotesingle{}}\NormalTok{] }\OperatorTok{\textless{}=} \DecValTok{30}\NormalTok{))][}\StringTok{\textquotesingle{}Prcp(mm)\textquotesingle{}}\NormalTok{].}\BuiltInTok{sum}\NormalTok{()}

    \CommentTok{\# Save weather data for the site}
\NormalTok{    weather\_filename }\OperatorTok{=}\NormalTok{ os.path.join(weather\_dir, }\SpecialStringTok{f"site\_}\SpecialCharTok{\{}\NormalTok{site}\SpecialCharTok{\}}\SpecialStringTok{\_weather\_data.txt"}\NormalTok{)}
\NormalTok{    df\_filtered.to\_csv(weather\_filename, sep}\OperatorTok{=}\StringTok{"}\CharTok{\textbackslash{}t}\StringTok{"}\NormalTok{, index}\OperatorTok{=}\VariableTok{False}\NormalTok{)}
    \BuiltInTok{print}\NormalTok{(}\SpecialStringTok{f"Weather data for site }\SpecialCharTok{\{}\NormalTok{site}\SpecialCharTok{\}}\SpecialStringTok{ saved as }\SpecialCharTok{\{}\NormalTok{weather\_filename}\SpecialCharTok{\}}\SpecialStringTok{"}\NormalTok{)}

    \CommentTok{\# Prepare weather data for AquaCrop}
\NormalTok{    wdf }\OperatorTok{=}\NormalTok{ prepare\_weather(weather\_filename)}

    \CommentTok{\# Define irrigation, soil, crop, and initial water content}
\NormalTok{    irr\_mngt }\OperatorTok{=}\NormalTok{ IrrigationManagement(irrigation\_method}\OperatorTok{=}\DecValTok{0}\NormalTok{)  }\CommentTok{\# Rain{-}fed}
\NormalTok{    soil }\OperatorTok{=}\NormalTok{ Soil(}\StringTok{\textquotesingle{}LoamySand\textquotesingle{}}\NormalTok{)}
\NormalTok{    crop }\OperatorTok{=}\NormalTok{ Crop(}\StringTok{\textquotesingle{}Wheat\textquotesingle{}}\NormalTok{,}
\NormalTok{                planting\_date}\OperatorTok{=}\StringTok{\textquotesingle{}04/01\textquotesingle{}}\NormalTok{,}
\NormalTok{                harvest\_date}\OperatorTok{=}\StringTok{\textquotesingle{}11/30\textquotesingle{}}\NormalTok{,}
\NormalTok{                CropType}\OperatorTok{=}\DecValTok{3}\NormalTok{,}
\NormalTok{                Tbase}\OperatorTok{=}\DecValTok{10}\NormalTok{,}
\NormalTok{                Tupp}\OperatorTok{=}\DecValTok{35}\NormalTok{,}
\NormalTok{                Zmax}\OperatorTok{=}\FloatTok{0.7}\NormalTok{,}
\NormalTok{                WP}\OperatorTok{=}\DecValTok{16}\NormalTok{,}
\NormalTok{                Tmin\_up}\OperatorTok{=}\DecValTok{8}\NormalTok{,}
\NormalTok{                Tmax\_lo}\OperatorTok{=}\DecValTok{40}\NormalTok{,}
\NormalTok{                exc}\OperatorTok{=}\DecValTok{50}\NormalTok{,}
\NormalTok{                CGC}\OperatorTok{=}\FloatTok{0.16764}\NormalTok{,}
\NormalTok{                CCx}\OperatorTok{=}\FloatTok{0.95}\NormalTok{,}
\NormalTok{                CDC}\OperatorTok{=}\FloatTok{0.13653}\NormalTok{,}
\NormalTok{                Kcb}\OperatorTok{=}\FloatTok{1.10}\NormalTok{,}
\NormalTok{                fshape\_r}\OperatorTok{=}\DecValTok{15}\NormalTok{,}
\NormalTok{                SxTopQ}\OperatorTok{=}\FloatTok{0.020}\NormalTok{,}
\NormalTok{                SxBotQ}\OperatorTok{=}\FloatTok{0.005}\NormalTok{,}
\NormalTok{                p\_up4}\OperatorTok{=}\FloatTok{0.8}\NormalTok{,}
\NormalTok{                p\_up2}\OperatorTok{=}\FloatTok{0.55}\NormalTok{,}
\NormalTok{                fshape\_w1}\OperatorTok{=}\DecValTok{4}
\NormalTok{               )}
\NormalTok{    initWC }\OperatorTok{=}\NormalTok{ InitialWaterContent(value}\OperatorTok{=}\NormalTok{[}\StringTok{\textquotesingle{}FC\textquotesingle{}}\NormalTok{])}

    \CommentTok{\# Create AquaCrop model instance}
\NormalTok{    model }\OperatorTok{=}\NormalTok{ AquaCropModel(sim\_start, sim\_end, wdf, soil, crop, initial\_water\_content}\OperatorTok{=}\NormalTok{initWC, irrigation\_management}\OperatorTok{=}\NormalTok{irr\_mngt)}

    \CommentTok{\# Run simulation}
    \ControlFlowTok{try}\NormalTok{:}
\NormalTok{        model.run\_model(till\_termination}\OperatorTok{=}\VariableTok{True}\NormalTok{)}
        \CommentTok{\# Create site{-}specific subdirectory}
\NormalTok{        site\_sim\_output\_dir }\OperatorTok{=}\NormalTok{ os.path.join(sim\_output\_dir, }\SpecialStringTok{f"site\_}\SpecialCharTok{\{}\NormalTok{site}\SpecialCharTok{\}}\SpecialStringTok{"}\NormalTok{)}
\NormalTok{        os.makedirs(site\_sim\_output\_dir, exist\_ok}\OperatorTok{=}\VariableTok{True}\NormalTok{)}\OperatorTok{;}

        \CommentTok{\# Save site results}
\NormalTok{        results\_df }\OperatorTok{=}\NormalTok{ model.\_outputs.final\_stats}
\NormalTok{        results\_df[}\StringTok{\textquotesingle{}Total\_Precipitation(mm)\textquotesingle{}}\NormalTok{] }\OperatorTok{=}\NormalTok{ total\_precip}
\NormalTok{        site\_results\_file }\OperatorTok{=}\NormalTok{ os.path.join(site\_sim\_output\_dir, }\SpecialStringTok{f"site\_}\SpecialCharTok{\{}\NormalTok{site}\SpecialCharTok{\}}\SpecialStringTok{\_simulation\_results.csv"}\NormalTok{)}
\NormalTok{        results\_df.to\_csv(site\_results\_file, index}\OperatorTok{=}\VariableTok{False}\NormalTok{)}
        \BuiltInTok{print}\NormalTok{(}\SpecialStringTok{f"Simulation results for site }\SpecialCharTok{\{}\NormalTok{site}\SpecialCharTok{\}}\SpecialStringTok{ saved in }\SpecialCharTok{\{}\NormalTok{site\_sim\_output\_dir}\SpecialCharTok{\}}\SpecialStringTok{"}\NormalTok{)}

        \CommentTok{\# Append to all results for later merging}
\NormalTok{        results\_df[}\StringTok{\textquotesingle{}Site\textquotesingle{}}\NormalTok{] }\OperatorTok{=}\NormalTok{ site}
\NormalTok{        all\_simulation\_results.append(results\_df)}

    \ControlFlowTok{except} \PreprocessorTok{Exception} \ImportTok{as}\NormalTok{ e:}
        \BuiltInTok{print}\NormalTok{(}\SpecialStringTok{f"Error during simulation for site }\SpecialCharTok{\{}\NormalTok{site}\SpecialCharTok{\}}\SpecialStringTok{: }\SpecialCharTok{\{}\NormalTok{e}\SpecialCharTok{\}}\SpecialStringTok{"}\NormalTok{)}
\end{Highlighting}
\end{Shaded}

\begin{Shaded}
\begin{Highlighting}[]
\CommentTok{\# {-}{-}{-} Merge all results {-}{-}{-}}
\ControlFlowTok{if}\NormalTok{ all\_simulation\_results:}
\NormalTok{    merged\_results }\OperatorTok{=}\NormalTok{ pd.concat(all\_simulation\_results, ignore\_index}\OperatorTok{=}\VariableTok{True}\NormalTok{)}
\NormalTok{    merged\_results.to\_csv(merged\_output\_file, index}\OperatorTok{=}\VariableTok{False}\NormalTok{)}
    \BuiltInTok{print}\NormalTok{(}\SpecialStringTok{f"All site simulations merged and saved as }\SpecialCharTok{\{}\NormalTok{merged\_output\_file}\SpecialCharTok{\}}\SpecialStringTok{"}\NormalTok{)}
\ControlFlowTok{else}\NormalTok{:}
    \BuiltInTok{print}\NormalTok{(}\StringTok{"No simulation results to merge."}\NormalTok{)}
\end{Highlighting}
\end{Shaded}

\textbf{Simulation Configuration 2}

Our second simulation configuration aims to simulate the crop production
function under an irrigation production system, simulating the crop
yield--water relationship under irrigation conditions.

Here, we keep the code the same as for rainfed conditions, except that
we change the irrigation management to:

irr\_mngt = IrrigationManagement(irrigation\_method=4, NetIrrSMT=70)

With these simulation results, we also obtained data on the net
irrigation amount at each site required to achieve the maximum potential
yield.

The same specifications are then applied for the years 2019, 2020, 2021,
2022, and 2023.

The same workflow is also applied for canola and potato, with crop
parameters customized for each crop based on published and calibrated
values.

For potato, we do not simulate rain-fed conditions, as growing potatoes
under rain-fed conditions in Saskatchewan is rare.

\textbf{Simulation Configuration 3}

In the third simulation configuration, we aim to generate the data
needed to estimate the marginal value of irrigation water. The marginal
value of irrigation water is calculated as:

\begin{equation}
    MV_{jt}=\sum_{j=1}^{n} \frac{\triangle NR^{IR_q}_{itj}}{{\triangle IRWQ_{it}}}
    \label{equation2}
\end{equation}

\noindent Where

\noindent \(NR^{IR_q}_{itj} = Y^{IR_q}_{itj}*P_{it} - C_{it}^{IR}\)

Once you understand the basics of AquaCrop-OSPy, we can move to a more
advanced application---optimizing irrigation strategies using real
climate data for wheat (and later canola and potato). In our project
around Lake Diefenbaker, Saskatchewan, we simulate crop yields under
both rain-fed and irrigation conditions and estimate how much irrigation
is needed to reach maximum potential yield. This section explains how we
use AquaCrop-OSPy together with an optimization routine to find the best
irrigation management strategy for each field given a limit on total
irrigation supply.

In earlier simulations, we assumed a fixed irrigation rule:

``Apply irrigation whenever soil moisture falls below a predefined
threshold.''

However, when water is limited, this rule does not necessarily maximize
yield. Instead, we need to find the best allocation of a fixed seasonal
irrigation amount. For this, we optimize soil-moisture thresholds (SMTs)
across the four major crop growth stages.

How this works in simulation framework;

\begin{enumerate}
\def\labelenumi{\arabic{enumi}.}
\tightlist
\item
  We choose maximum irrigation allowance per season (i.e.~20mm, 50mm ,
  100mm, etc)
\item
  Optimizer search for the best 4 soil moisture tresholds (SMTs)

  \begin{itemize}
  \tightlist
  \item
    one threshold for each crop growth stage: (emergence, canopy
    growth,maximum canopy and senescence)
  \end{itemize}
\item
  Whenever soil moisture goes below a threshold, irrigation is applied.
\item
  Why SMTs? Because :
\end{enumerate}

\begin{itemize}
\tightlist
\item
  Crops need more water during certain growth stages.
\item
  Stress at key stages (e.g., flowering) reduces yield more than at
  others.
\item
  Using different thresholds allows dynamic irrigation scheduling.
\end{itemize}

The Optimizer (scipy.optimize) specifically;

\begin{itemize}
\tightlist
\item
  Try many different SMT combinations
\item
  Simulate AquaCrop for each combination
\item
  Measure the yield
\item
  Find the SMT combination that gives the highest yield under the water
  constraint
\end{itemize}

\begin{enumerate}
\def\labelenumi{\arabic{enumi}.}
\setcounter{enumi}{4}
\tightlist
\item
  Repeating Optimization for Different Water Constraints
\end{enumerate}

For example: - Max irrigation = 50 mm → find best SMTs → Highest yiels
under 50mm irrigation availabiltiy - Max irrigation = 100 mm → find best
SMTs → Highest yiels under 100mm irrigation availabiltiy - Max
irrigation = 200 mm → find best SMTs → Highest yiels under 200mm
irrigation availabiltiy

\textbf{Step 1}

Create climate data files with ET₀ --- this step is similar to creating
the climate data files in the above simulation configurations.

\begin{Shaded}
\begin{Highlighting}[]
\ImportTok{import}\NormalTok{ pandas }\ImportTok{as}\NormalTok{ pd}
\ImportTok{import}\NormalTok{ numpy }\ImportTok{as}\NormalTok{ np}

\CommentTok{\# Load the downloaded Daymet data}
\NormalTok{climate\_data }\OperatorTok{=}\NormalTok{ pd.read\_csv(}\StringTok{\textquotesingle{}/Users/tharakajayalath/Library/CloudStorage/OneDrive{-}UniversityofSaskatchewan/Chapter II{-}IrrigationValue/Chapter{-}II/AquaCropOPSyData/ClimateData/weather\_data\_Aqua.csv\textquotesingle{}}\NormalTok{, on\_bad\_lines}\OperatorTok{=}\StringTok{\textquotesingle{}skip\textquotesingle{}}\NormalTok{)}

\CommentTok{\# Check if \textquotesingle{}Date\textquotesingle{} column exists}
\ControlFlowTok{if} \StringTok{\textquotesingle{}Date\textquotesingle{}} \KeywordTok{in}\NormalTok{ climate\_data.columns:}
    \CommentTok{\# Convert \textquotesingle{}Date\textquotesingle{} column to datetime}
\NormalTok{    climate\_data[}\StringTok{\textquotesingle{}Date\textquotesingle{}}\NormalTok{] }\OperatorTok{=}\NormalTok{ pd.to\_datetime(climate\_data[}\StringTok{\textquotesingle{}Date\textquotesingle{}}\NormalTok{], errors}\OperatorTok{=}\StringTok{\textquotesingle{}coerce\textquotesingle{}}\NormalTok{)}

    \CommentTok{\# Remove rows with missing Date values}
\NormalTok{    climate\_data }\OperatorTok{=}\NormalTok{ climate\_data.dropna(subset}\OperatorTok{=}\NormalTok{[}\StringTok{\textquotesingle{}Date\textquotesingle{}}\NormalTok{])}

    \CommentTok{\# Extract Day, Month, and Year from the Date column}
\NormalTok{    climate\_data[}\StringTok{\textquotesingle{}Day\textquotesingle{}}\NormalTok{] }\OperatorTok{=}\NormalTok{ climate\_data[}\StringTok{\textquotesingle{}Date\textquotesingle{}}\NormalTok{].dt.day}
\NormalTok{    climate\_data[}\StringTok{\textquotesingle{}Month\textquotesingle{}}\NormalTok{] }\OperatorTok{=}\NormalTok{ climate\_data[}\StringTok{\textquotesingle{}Date\textquotesingle{}}\NormalTok{].dt.month}
\NormalTok{    climate\_data[}\StringTok{\textquotesingle{}Year\textquotesingle{}}\NormalTok{] }\OperatorTok{=}\NormalTok{ climate\_data[}\StringTok{\textquotesingle{}Date\textquotesingle{}}\NormalTok{].dt.year}
\NormalTok{    e\_a }\OperatorTok{=}\NormalTok{ climate\_data[}\StringTok{\textquotesingle{}e\_a\textquotesingle{}}\NormalTok{]  }\CommentTok{\# actual vapor pressure (kpa)}

    \CommentTok{\# Calculate necessary climate variables}
\NormalTok{    T\_max }\OperatorTok{=}\NormalTok{ climate\_data[}\StringTok{\textquotesingle{}MaxTemp\textquotesingle{}}\NormalTok{]  }\CommentTok{\# Maximum temperature (°C)}
\NormalTok{    T\_min }\OperatorTok{=}\NormalTok{ climate\_data[}\StringTok{\textquotesingle{}MinTemp\textquotesingle{}}\NormalTok{]  }\CommentTok{\# Minimum temperature (°C)}
\NormalTok{    precipitation }\OperatorTok{=}\NormalTok{ climate\_data[}\StringTok{\textquotesingle{}Precipitation\textquotesingle{}}\NormalTok{]  }\CommentTok{\# Precipitation (mm)}
\NormalTok{    solar\_radiation }\OperatorTok{=}\NormalTok{ climate\_data[}\StringTok{\textquotesingle{}R\_n\textquotesingle{}}\NormalTok{]  }\CommentTok{\# Solar radiation (MJ/m²/day)}

    \CommentTok{\# Calculate mean temperature}
\NormalTok{    T\_mean }\OperatorTok{=}\NormalTok{ (T\_max }\OperatorTok{+}\NormalTok{ T\_min) }\OperatorTok{/} \DecValTok{2}

    \CommentTok{\# Calculate saturation vapor pressure (e\_s) in kPa}
\NormalTok{    e\_s }\OperatorTok{=} \FloatTok{0.6108} \OperatorTok{*}\NormalTok{ np.exp((}\FloatTok{17.27} \OperatorTok{*}\NormalTok{ T\_mean) }\OperatorTok{/}\NormalTok{ (T\_mean }\OperatorTok{+} \FloatTok{237.3}\NormalTok{))}

    \CommentTok{\# Assuming relative humidity (RH) is available or use an average value}
    \CommentTok{\#RH = 62  \# Example average RH of 50\%}
    \CommentTok{\#e\_a = (RH / 100) * e\_s  \# Actual vapor pressure in kPa}

    \CommentTok{\# Define other constants}
\NormalTok{    R\_n }\OperatorTok{=}\NormalTok{ solar\_radiation  }\CommentTok{\# Net radiation (MJ/m²/day)}
\NormalTok{    G }\OperatorTok{=} \DecValTok{0}  \CommentTok{\# Soil heat flux density (MJ/m²/day), often negligible}
\NormalTok{    gamma }\OperatorTok{=} \FloatTok{0.066}  \CommentTok{\# Psychrometric constant (kPa/°C)}
\NormalTok{    u }\OperatorTok{=} \DecValTok{2}  \CommentTok{\# Average wind speed in m/s (if available, else use an estimated value)}

    \CommentTok{\# Calculate slope of the saturation vapor pressure curve (Δ)}
\NormalTok{    delta }\OperatorTok{=}\NormalTok{ (}\DecValTok{4098} \OperatorTok{*}\NormalTok{ e\_s) }\OperatorTok{/}\NormalTok{ ((T\_mean }\OperatorTok{+} \FloatTok{237.3}\NormalTok{) }\OperatorTok{**} \DecValTok{2}\NormalTok{)}

    \CommentTok{\# Calculate reference ET (ET₀) using the Penman{-}Monteith equation}
\NormalTok{    ReferenceET }\OperatorTok{=}\NormalTok{ (}\FloatTok{0.408} \OperatorTok{*}\NormalTok{ delta }\OperatorTok{*}\NormalTok{ (R\_n }\OperatorTok{{-}}\NormalTok{ G) }\OperatorTok{+}\NormalTok{ gamma }\OperatorTok{*}\NormalTok{ (}\DecValTok{900} \OperatorTok{/}\NormalTok{ (T\_mean }\OperatorTok{+} \DecValTok{273}\NormalTok{)) }\OperatorTok{*}\NormalTok{ u }\OperatorTok{*}\NormalTok{ (e\_s }\OperatorTok{{-}}\NormalTok{ e\_a)) }\OperatorTok{/}\NormalTok{ (delta }\OperatorTok{+}\NormalTok{ gamma }\OperatorTok{*}\NormalTok{ (}\DecValTok{1} \OperatorTok{+} \FloatTok{0.34} \OperatorTok{*}\NormalTok{ u))}

    \CommentTok{\# Add ET₀ to the DataFrame}
\NormalTok{    climate\_data[}\StringTok{\textquotesingle{}ReferenceET\textquotesingle{}}\NormalTok{] }\OperatorTok{=}\NormalTok{ ReferenceET}

    \CommentTok{\# Reorder the columns to bring Day, Month, Year first}
\NormalTok{    climate\_data }\OperatorTok{=}\NormalTok{ climate\_data[[}\StringTok{\textquotesingle{}Day\textquotesingle{}}\NormalTok{, }\StringTok{\textquotesingle{}Month\textquotesingle{}}\NormalTok{, }\StringTok{\textquotesingle{}Year\textquotesingle{}}\NormalTok{] }\OperatorTok{+}\NormalTok{ [col }\ControlFlowTok{for}\NormalTok{ col }\KeywordTok{in}\NormalTok{ climate\_data.columns }\ControlFlowTok{if}\NormalTok{ col }\KeywordTok{not} \KeywordTok{in}\NormalTok{ [}\StringTok{\textquotesingle{}Day\textquotesingle{}}\NormalTok{, }\StringTok{\textquotesingle{}Month\textquotesingle{}}\NormalTok{, }\StringTok{\textquotesingle{}Year\textquotesingle{}}\NormalTok{]]]}

    \CommentTok{\# Save the DataFrame to a text file (CSV format)}
\NormalTok{    climate\_data.to\_csv(}\StringTok{\textquotesingle{}daymet\_data\_with\_et0.csv\textquotesingle{}}\NormalTok{, index}\OperatorTok{=}\VariableTok{False}\NormalTok{)  }\CommentTok{\# Save as a tab{-}separated text file}

    \BuiltInTok{print}\NormalTok{(}\StringTok{"Reference ET₀ calculated and saved to \textquotesingle{}daymet\_data\_with\_et0.csv\textquotesingle{}."}\NormalTok{)}
\ControlFlowTok{else}\NormalTok{:}
    \BuiltInTok{print}\NormalTok{(}\StringTok{"The \textquotesingle{}Date\textquotesingle{} column does not exist in the dataset."}\NormalTok{)}
\end{Highlighting}
\end{Shaded}

\begin{verbatim}
Reference ET₀ calculated and saved to 'daymet_data_with_et0.csv'.
\end{verbatim}

\textbf{Step 2}

In contrast to simulating all 394 sites at once, running optimization
algorithms to find the optimal soil-moisture thresholds (SMTs) and the
best yield under given irrigation levels for every site simultaneously
can be very memory-intensive. To handle this efficiently, we create
separate climate data files for each site. This allows us to run
simulations batch-wise, improving memory allocation and computational
efficiency.

Here's what the workflow in the code does:

\begin{enumerate}
\def\labelenumi{\arabic{enumi}.}
\tightlist
\item
  Start with the master CSV:
\end{enumerate}

We begin with a CSV file (daymet\_data\_with\_et0.csv) that contains
daily climate data for all sites. Each row includes a site identifier
(site) along with day, month, year, and climate variables.

\begin{enumerate}
\def\labelenumi{\arabic{enumi}.}
\setcounter{enumi}{1}
\tightlist
\item
  Split into site-specific climate files:
\end{enumerate}

The code loops through each unique site and writes a separate text file
for it. For example, for site 1, a file ClimateData/site\_1.txt is
created, for site 2, ClimateData/site\_2.txt, and so on. Each file
contains only the climate data for that particular site.

\begin{enumerate}
\def\labelenumi{\arabic{enumi}.}
\setcounter{enumi}{2}
\tightlist
\item
  Prepare weather data for AquaCrop:
\end{enumerate}

After creating site-specific files, we use the prepare\_weather()
function on each file. This converts the raw climate data into a format
that AquaCrop-OSPy can directly use for simulation.

\begin{enumerate}
\def\labelenumi{\arabic{enumi}.}
\setcounter{enumi}{3}
\tightlist
\item
  Simulate per site:
\end{enumerate}

With the prepared weather data, we can run AquaCrop simulations
independently for each site. This allows us to optimize irrigation
strategies and predict crop yields tailored to the specific conditions
at each location.

By processing the data site by site, we reduce memory load and make it
feasible to run complex optimization algorithms across hundreds of sites
efficiently.

\begin{Shaded}
\begin{Highlighting}[]
\ImportTok{import}\NormalTok{ os}
\ImportTok{import}\NormalTok{ pandas }\ImportTok{as}\NormalTok{ pd}

\CommentTok{\# {-}{-}{-} Load data {-}{-}{-}}
\NormalTok{input\_file }\OperatorTok{=} \StringTok{"./AquaCropReadme/daymet\_data\_with\_et0.csv"}
\NormalTok{df }\OperatorTok{=}\NormalTok{ pd.read\_csv(input\_file)}

\CommentTok{\# {-}{-}{-} Unique sites {-}{-}{-}}
\NormalTok{unique\_sites }\OperatorTok{=}\NormalTok{ df[}\StringTok{\textquotesingle{}site\textquotesingle{}}\NormalTok{].unique()}

\CommentTok{\# {-}{-}{-} Directory for weather files (same as earlier setup) {-}{-}{-}}
\NormalTok{cwd }\OperatorTok{=}\NormalTok{ os.getcwd()}
\NormalTok{weather\_dir }\OperatorTok{=}\NormalTok{ os.path.join(cwd, }\StringTok{"ClimateData"}\NormalTok{)}
\NormalTok{os.makedirs(weather\_dir, exist\_ok}\OperatorTok{=}\VariableTok{True}\NormalTok{)}

\CommentTok{\# {-}{-}{-} Columns to select and rename {-}{-}{-}}
\NormalTok{columns\_to\_select }\OperatorTok{=}\NormalTok{ [}\StringTok{\textquotesingle{}Day\textquotesingle{}}\NormalTok{, }\StringTok{\textquotesingle{}Month\textquotesingle{}}\NormalTok{, }\StringTok{\textquotesingle{}Year\textquotesingle{}}\NormalTok{, }\StringTok{\textquotesingle{}MinTemp\textquotesingle{}}\NormalTok{, }\StringTok{\textquotesingle{}MaxTemp\textquotesingle{}}\NormalTok{, }\StringTok{\textquotesingle{}Precipitation\textquotesingle{}}\NormalTok{, }\StringTok{\textquotesingle{}ReferenceET\textquotesingle{}}\NormalTok{]}
\NormalTok{columns\_to\_rename }\OperatorTok{=}\NormalTok{ \{}
    \StringTok{\textquotesingle{}MinTemp\textquotesingle{}}\NormalTok{: }\StringTok{\textquotesingle{}Tmin(c)\textquotesingle{}}\NormalTok{,}
    \StringTok{\textquotesingle{}MaxTemp\textquotesingle{}}\NormalTok{: }\StringTok{\textquotesingle{}Tmax(c)\textquotesingle{}}\NormalTok{,}
    \StringTok{\textquotesingle{}Precipitation\textquotesingle{}}\NormalTok{: }\StringTok{\textquotesingle{}Prcp(mm)\textquotesingle{}}\NormalTok{,}
    \StringTok{\textquotesingle{}ReferenceET\textquotesingle{}}\NormalTok{: }\StringTok{\textquotesingle{}Et0(mm)\textquotesingle{}}
\NormalTok{\}}

\CommentTok{\# {-}{-}{-} Loop through each site and save weather files {-}{-}{-}}
\ControlFlowTok{for}\NormalTok{ site }\KeywordTok{in}\NormalTok{ unique\_sites:}
\NormalTok{    df\_filtered }\OperatorTok{=}\NormalTok{ (}
\NormalTok{        df[df[}\StringTok{\textquotesingle{}site\textquotesingle{}}\NormalTok{] }\OperatorTok{==}\NormalTok{ site][columns\_to\_select]}
\NormalTok{        .rename(columns}\OperatorTok{=}\NormalTok{columns\_to\_rename)}
\NormalTok{    )}
    
    \CommentTok{\# Weather file path}
\NormalTok{    weather\_filename }\OperatorTok{=}\NormalTok{ os.path.join(weather\_dir, }\SpecialStringTok{f"site\_}\SpecialCharTok{\{}\NormalTok{site}\SpecialCharTok{\}}\SpecialStringTok{.txt"}\NormalTok{)}
    
    \CommentTok{\# Save weather file}
\NormalTok{    df\_filtered.to\_csv(weather\_filename, sep}\OperatorTok{=}\StringTok{"}\CharTok{\textbackslash{}t}\StringTok{"}\NormalTok{, index}\OperatorTok{=}\VariableTok{False}\NormalTok{)}
    \BuiltInTok{print}\NormalTok{(}\SpecialStringTok{f"Weather data saved for site }\SpecialCharTok{\{}\NormalTok{site}\SpecialCharTok{\}}\SpecialStringTok{."}\NormalTok{)  }\CommentTok{\# Remove if you want full silence}
\end{Highlighting}
\end{Shaded}

\textbf{Step 3}

Create whether data file for each site.

\begin{Shaded}
\begin{Highlighting}[]
\ImportTok{import}\NormalTok{ os}
\ImportTok{from}\NormalTok{ aquacrop.utils.prepare\_weather }\ImportTok{import}\NormalTok{ prepare\_weather}

\CommentTok{\# Base working directory (same as earlier)}
\NormalTok{cwd }\OperatorTok{=}\NormalTok{ os.getcwd()}

\CommentTok{\# Input directory containing the raw climate text files}
\NormalTok{input\_dir }\OperatorTok{=}\NormalTok{ os.path.join(cwd, }\StringTok{"ClimateData"}\NormalTok{)}

\CommentTok{\# Output directory for prepared WDF files}
\NormalTok{prepared\_wdf\_dir }\OperatorTok{=}\NormalTok{ os.path.join(cwd, }\StringTok{"PreparedWDF"}\NormalTok{)}
\NormalTok{os.makedirs(prepared\_wdf\_dir, exist\_ok}\OperatorTok{=}\VariableTok{True}\NormalTok{)}

\CommentTok{\# Number of sites}
\NormalTok{num\_sites }\OperatorTok{=} \DecValTok{397}  \CommentTok{\# adjust if needed}

\CommentTok{\# Dynamically create file names (site\_1.txt ... site\_397.txt)}
\NormalTok{climate\_files }\OperatorTok{=}\NormalTok{ [}\SpecialStringTok{f"site\_}\SpecialCharTok{\{}\NormalTok{i}\SpecialCharTok{\}}\SpecialStringTok{.txt"} \ControlFlowTok{for}\NormalTok{ i }\KeywordTok{in} \BuiltInTok{range}\NormalTok{(}\DecValTok{1}\NormalTok{, num\_sites }\OperatorTok{+} \DecValTok{1}\NormalTok{)]}

\CommentTok{\# Loop and process each climate file}
\ControlFlowTok{for}\NormalTok{ climate\_file }\KeywordTok{in}\NormalTok{ climate\_files:}
\NormalTok{    file\_path }\OperatorTok{=}\NormalTok{ os.path.join(input\_dir, climate\_file)}

    \ControlFlowTok{if}\NormalTok{ os.path.exists(file\_path):}
        \ControlFlowTok{try}\NormalTok{:}
            \CommentTok{\# Prepare weather data for AquaCrop}
\NormalTok{            wdf }\OperatorTok{=}\NormalTok{ prepare\_weather(file\_path)}

            \CommentTok{\# Save the prepared weather file}
\NormalTok{            prepared\_path }\OperatorTok{=}\NormalTok{ os.path.join(prepared\_wdf\_dir, }\SpecialStringTok{f"prepared\_}\SpecialCharTok{\{}\NormalTok{climate\_file}\SpecialCharTok{\}}\SpecialStringTok{"}\NormalTok{)}
\NormalTok{            wdf.to\_csv(prepared\_path, index}\OperatorTok{=}\VariableTok{False}\NormalTok{)}

            \BuiltInTok{print}\NormalTok{(}\SpecialStringTok{f"Weather data prepared and saved: }\SpecialCharTok{\{}\NormalTok{climate\_file}\SpecialCharTok{\}}\SpecialStringTok{"}\NormalTok{)}

        \ControlFlowTok{except} \PreprocessorTok{Exception} \ImportTok{as}\NormalTok{ e:}
            \BuiltInTok{print}\NormalTok{(}\SpecialStringTok{f"Error processing }\SpecialCharTok{\{}\NormalTok{climate\_file}\SpecialCharTok{\}}\SpecialStringTok{: }\SpecialCharTok{\{}\NormalTok{e}\SpecialCharTok{\}}\SpecialStringTok{"}\NormalTok{)}

    \ControlFlowTok{else}\NormalTok{:}
        \BuiltInTok{print}\NormalTok{(}\SpecialStringTok{f"Missing file: }\SpecialCharTok{\{}\NormalTok{file\_path}\SpecialCharTok{\}}\SpecialStringTok{"}\NormalTok{)}
\end{Highlighting}
\end{Shaded}

\textbf{Step 3}

\textbf{Understanding and Defining the Core Functions}

\textbf{1. run\_model()}

Runs the AquaCrop model using:

\begin{itemize}
\item
  a set of soil-moisture thresholds (SMTs)
\item
  a maximum seasonal irrigation limit
\item
  crop and soil parameters
\item
  weather data
\item
  simulation start and end dates
\end{itemize}

It returns simulation outputs, including yield and irrigation applied.

\textbf{2. evaluate()}

Takes SMTs and:

\begin{itemize}
\item
  runs the model
\item
  extracts mean yield and total irrigation
\end{itemize}

If test=False, it returns -yield, because: The optimizer minimizes the
objective, so minimizing -yield = maximizing yield.

If test=True, it returns the actual numbers: yield, irrigation, reward.

\textbf{3. get\_starting\_point()}

Before optimization, we generate many random SMT combinations and
evaluate their performance.

The best-performing one becomes the starting point for the local
optimization.

This avoids getting stuck in poor solutions.

\textbf{4. optimize()}

This is the main optimization function. It:

\begin{itemize}
\item
  Gets a good initial SMT guess using get\_starting\_point()
\item
  Uses scipy.optimize.fmin() to refine the SMTs and maximize yield
\item
  Returns the optimal SMTs for the given irrigation limit and weather
  data
\end{itemize}

\begin{Shaded}
\begin{Highlighting}[]
\ImportTok{import}\NormalTok{ os}
\ImportTok{import}\NormalTok{ numpy }\ImportTok{as}\NormalTok{ np}
\ImportTok{import}\NormalTok{ pandas }\ImportTok{as}\NormalTok{ pd}
\ImportTok{from}\NormalTok{ aquacrop }\ImportTok{import}\NormalTok{ AquaCropModel, Soil, Crop, InitialWaterContent, IrrigationManagement}
\ImportTok{from}\NormalTok{ scipy.optimize }\ImportTok{import}\NormalTok{ fmin}


\CommentTok{\# {-}{-}{-}{-}{-}{-}{-}{-}{-}{-}{-}{-}{-}{-}{-}{-}{-}{-}{-}{-}{-}{-}{-}{-}{-}{-}{-}{-}{-}{-}{-}{-}{-}{-}{-}{-}{-}{-}{-}{-}{-}{-}{-}{-}}
\CommentTok{\# 1. Define run\_model() function}
\CommentTok{\# {-}{-}{-}{-}{-}{-}{-}{-}{-}{-}{-}{-}{-}{-}{-}{-}{-}{-}{-}{-}{-}{-}{-}{-}{-}{-}{-}{-}{-}{-}{-}{-}{-}{-}{-}{-}{-}{-}{-}{-}{-}{-}{-}{-}}

\KeywordTok{def}\NormalTok{ run\_model(smts, max\_irr\_season, year1, year2, wdf):}
\NormalTok{    wheat }\OperatorTok{=}\NormalTok{ Crop(}
        \StringTok{\textquotesingle{}Wheat\textquotesingle{}}\NormalTok{,}
\NormalTok{        planting\_date}\OperatorTok{=}\StringTok{\textquotesingle{}05/01\textquotesingle{}}\NormalTok{,}
\NormalTok{        harvest\_date}\OperatorTok{=}\StringTok{\textquotesingle{}10/30\textquotesingle{}}\NormalTok{,}
\NormalTok{        CropType}\OperatorTok{=}\DecValTok{3}\NormalTok{,}
\NormalTok{        Tbase}\OperatorTok{=}\DecValTok{10}\NormalTok{,}
\NormalTok{        Tupp}\OperatorTok{=}\DecValTok{35}\NormalTok{,}
\NormalTok{        Zmax}\OperatorTok{=}\FloatTok{0.7}\NormalTok{,}
\NormalTok{        WP}\OperatorTok{=}\DecValTok{16}\NormalTok{,}
\NormalTok{        Tmin\_up}\OperatorTok{=}\DecValTok{8}\NormalTok{,}
\NormalTok{        Tmax\_lo}\OperatorTok{=}\DecValTok{40}\NormalTok{,}
\NormalTok{        exc}\OperatorTok{=}\DecValTok{50}\NormalTok{,}
\NormalTok{        CGC}\OperatorTok{=}\FloatTok{0.16764}\NormalTok{,}
\NormalTok{        CCx}\OperatorTok{=}\FloatTok{0.95}\NormalTok{,}
\NormalTok{        CDC}\OperatorTok{=}\FloatTok{0.13653}\NormalTok{,}
\NormalTok{        Kcb}\OperatorTok{=}\FloatTok{1.10}\NormalTok{,}
\NormalTok{        fshape\_r}\OperatorTok{=}\DecValTok{15}\NormalTok{,}
\NormalTok{        SxTopQ}\OperatorTok{=}\FloatTok{0.020}\NormalTok{,}
\NormalTok{        SxBotQ}\OperatorTok{=}\FloatTok{0.005}\NormalTok{,}
\NormalTok{        p\_up4}\OperatorTok{=}\FloatTok{0.8}\NormalTok{,}
\NormalTok{        p\_up2}\OperatorTok{=}\FloatTok{0.55}\NormalTok{,}
\NormalTok{        fshape\_w1}\OperatorTok{=}\DecValTok{4}
\NormalTok{    )}

\NormalTok{    soil }\OperatorTok{=}\NormalTok{ Soil(}\StringTok{\textquotesingle{}LoamySand\textquotesingle{}}\NormalTok{)}
\NormalTok{    init\_wc }\OperatorTok{=}\NormalTok{ InitialWaterContent(wc\_type}\OperatorTok{=}\StringTok{\textquotesingle{}Pct\textquotesingle{}}\NormalTok{, value}\OperatorTok{=}\NormalTok{[}\DecValTok{70}\NormalTok{])}

\NormalTok{    irrmngt }\OperatorTok{=}\NormalTok{ IrrigationManagement(}
\NormalTok{        irrigation\_method}\OperatorTok{=}\DecValTok{1}\NormalTok{,}
\NormalTok{        SMT}\OperatorTok{=}\NormalTok{smts,}
\NormalTok{        MaxIrrSeason}\OperatorTok{=}\NormalTok{max\_irr\_season}
\NormalTok{    )}

\NormalTok{    model }\OperatorTok{=}\NormalTok{ AquaCropModel(}
        \SpecialStringTok{f"}\SpecialCharTok{\{}\NormalTok{year1}\SpecialCharTok{\}}\SpecialStringTok{/05/01"}\NormalTok{,}
        \SpecialStringTok{f"}\SpecialCharTok{\{}\NormalTok{year2}\SpecialCharTok{\}}\SpecialStringTok{/10/31"}\NormalTok{,}
\NormalTok{        wdf,}
\NormalTok{        soil,}
\NormalTok{        wheat,}
\NormalTok{        irrigation\_management}\OperatorTok{=}\NormalTok{irrmngt,}
\NormalTok{        initial\_water\_content}\OperatorTok{=}\NormalTok{init\_wc}
\NormalTok{    )}

\NormalTok{    model.run\_model(till\_termination}\OperatorTok{=}\VariableTok{True}\NormalTok{)}
    \ControlFlowTok{return}\NormalTok{ model.get\_simulation\_results()}

\CommentTok{\# {-}{-}{-}{-}{-}{-}{-}{-}{-}{-}{-}{-}{-}{-}{-}{-}{-}{-}{-}{-}{-}{-}{-}{-}{-}{-}{-}{-}{-}{-}{-}{-}{-}{-}{-}{-}{-}{-}{-}{-}{-}{-}{-}{-}}
\CommentTok{\# 2. Define evaluate() function}
\CommentTok{\# {-}{-}{-}{-}{-}{-}{-}{-}{-}{-}{-}{-}{-}{-}{-}{-}{-}{-}{-}{-}{-}{-}{-}{-}{-}{-}{-}{-}{-}{-}{-}{-}{-}{-}{-}{-}{-}{-}{-}{-}{-}{-}{-}{-}}
\KeywordTok{def}\NormalTok{ evaluate(smts, max\_irr\_season, wdf, test}\OperatorTok{=}\VariableTok{False}\NormalTok{):}
\NormalTok{    out }\OperatorTok{=}\NormalTok{ run\_model(smts, max\_irr\_season, }\DecValTok{2018}\NormalTok{, }\DecValTok{2018}\NormalTok{, wdf)}

\NormalTok{    yld }\OperatorTok{=}\NormalTok{ out[}\StringTok{\textquotesingle{}Fresh yield (tonne/ha)\textquotesingle{}}\NormalTok{].mean()}
\NormalTok{    tirr }\OperatorTok{=}\NormalTok{ out[}\StringTok{\textquotesingle{}Seasonal irrigation (mm)\textquotesingle{}}\NormalTok{].mean()}

\NormalTok{    reward }\OperatorTok{=}\NormalTok{ yld}

    \ControlFlowTok{if}\NormalTok{ test:}
        \ControlFlowTok{return}\NormalTok{ yld, tirr, reward}
    \ControlFlowTok{else}\NormalTok{:}
        \ControlFlowTok{return} \OperatorTok{{-}}\NormalTok{reward}

\CommentTok{\# {-}{-}{-}{-}{-}{-}{-}{-}{-}{-}{-}{-}{-}{-}{-}{-}{-}{-}{-}{-}{-}{-}{-}{-}{-}{-}{-}{-}{-}{-}{-}{-}{-}{-}{-}{-}{-}{-}{-}{-}{-}{-}{-}{-}}
\CommentTok{\# 3. Define get\_starting\_point() function}
\CommentTok{\# {-}{-}{-}{-}{-}{-}{-}{-}{-}{-}{-}{-}{-}{-}{-}{-}{-}{-}{-}{-}{-}{-}{-}{-}{-}{-}{-}{-}{-}{-}{-}{-}{-}{-}{-}{-}{-}{-}{-}{-}{-}{-}{-}{-}}
\KeywordTok{def}\NormalTok{ get\_starting\_point(num\_smts, max\_irr\_season, num\_searches, wdf):}
\NormalTok{    x0list }\OperatorTok{=}\NormalTok{ np.random.rand(num\_searches, num\_smts) }\OperatorTok{*} \DecValTok{100} \CommentTok{\# Each SMT threshold is allowed to vary between 0\% and 100\% soil moisture.This defines the search space for irrigation thresholds.}
\NormalTok{    rlist }\OperatorTok{=}\NormalTok{ []}

    \ControlFlowTok{for}\NormalTok{ xtest }\KeywordTok{in}\NormalTok{ x0list:}
\NormalTok{        r }\OperatorTok{=}\NormalTok{ evaluate(xtest, max\_irr\_season, wdf)}
\NormalTok{        rlist.append(r)}

    \ControlFlowTok{return}\NormalTok{ x0list[np.argmin(rlist)]}

\CommentTok{\# num\_searches=100 means:}
\CommentTok{\# The algorithm will generate 100 random sets of SMT values}
\CommentTok{\# Evaluate all 100}
\CommentTok{\# Pick the best one (the one that gives the highest yield)}
\CommentTok{\# Use that best set as the starting point for the optimizer (fmin).}
\CommentTok{\# Why do we need 100 random starting points?}
\CommentTok{\# Optimizing irrigation thresholds is a nonlinear problem, so:}
\CommentTok{\# The objective function (yield) may have many local maxima.}
\CommentTok{\# If you start optimization from a bad initial guess, you might get stuck in a local optimum.}
\CommentTok{\# Trying 100 random initial guesses increases the chance of finding a good region of the search space.}

\CommentTok{\# {-}{-}{-}{-}{-}{-}{-}{-}{-}{-}{-}{-}{-}{-}{-}{-}{-}{-}{-}{-}{-}{-}{-}{-}{-}{-}{-}{-}{-}{-}{-}{-}{-}{-}{-}{-}{-}{-}{-}{-}{-}{-}{-}{-}}
\CommentTok{\# 4. Define optimize() function}
\CommentTok{\# {-}{-}{-}{-}{-}{-}{-}{-}{-}{-}{-}{-}{-}{-}{-}{-}{-}{-}{-}{-}{-}{-}{-}{-}{-}{-}{-}{-}{-}{-}{-}{-}{-}{-}{-}{-}{-}{-}{-}{-}{-}{-}{-}{-}}
\KeywordTok{def}\NormalTok{ optimize(num\_smts, max\_irr\_season, wdf, num\_searches}\OperatorTok{=}\DecValTok{100}\NormalTok{):}
\NormalTok{    x0 }\OperatorTok{=}\NormalTok{ get\_starting\_point(num\_smts, max\_irr\_season, num\_searches, wdf)}
\NormalTok{    res }\OperatorTok{=}\NormalTok{ fmin(evaluate, x0, disp}\OperatorTok{=}\DecValTok{0}\NormalTok{, args}\OperatorTok{=}\NormalTok{(max\_irr\_season, wdf))}
    \ControlFlowTok{return}\NormalTok{ res.squeeze()}
\end{Highlighting}
\end{Shaded}

In our set up with filed level simulationas and over 300 fields
following code optimize irrigation strategies for multiple sites and
irrigation levels in parallel, and save the results in a single CSV
file.

\textbf{1. Define irrigation levels}

max\_irr\_values =
{[}10,20,30,40,50,60,70,80,90,100,110,120,130,140,150,160,170,180,190,200{]}

These are maximum seasonal irrigation levels (mm) to test. The code will
optimize soil-moisture thresholds (SMTs) for each irrigation limit.

\textbf{2. Define site IDs}

site\_ids = range(1, 5)

Here we have over 300 sites and I have run the code for chunks of sites
as feeding higher number of sites into sigle lool ovewhem the memory

\textbf{3. Load prepared weather files}

wdfs = \{ site\_id:
prepare\_weather(f''./ClimateData/site\_\{site\_id\}.txt'') for site\_id
in site\_ids \}

\begin{itemize}
\tightlist
\item
  Loads site-specific weather files into a dictionary (wdfs) for quick
  access.
\item
  Each site has its own prepared weather data (wdf) for AquaCrop.
\end{itemize}

\textbf{4. Define the function to process each site + irrigation level}

def process\_site\_irr\_max(site\_id, max\_irr, wdf): smts = optimize(4,
max\_irr, wdf)\\
yld, tirr, \_ = evaluate(smts, max\_irr, wdf, True)

\begin{verbatim}
return {
    'Site_ID': site_id,
    'Max_Irrigation_mm': max_irr,
    'Optimal_SMTs': smts.tolist(),
    'Yield_tonne_per_ha': yld,
    'Total_Irrigation_mm': tirr
}

Inputs: site ID, max irrigation allowed, prepared weather data.
\end{verbatim}

Process:

\begin{itemize}
\item
  Calls optimize to find the best soil-moisture thresholds (SMTs) for
  the site and irrigation limit.
\item
  Calls evaluate with test=True to get yield and total irrigation for
  the optimized SMTs.
\item
  Output: a dictionary with site ID, irrigation level, optimal SMTs,
  yield, and irrigation applied.
\end{itemize}

\textbf{5. Run the optimization in parallel}

all\_results = Parallel(n\_jobs=15)(
delayed(process\_site\_irr\_max)(site\_id, max\_irr, wdfs{[}site\_id{]})
for max\_irr in max\_irr\_values for site\_id in site\_ids )

\begin{itemize}
\tightlist
\item
  Uses joblib to run multiple optimizations simultaneously.
\item
  n\_jobs=15 → runs 15 tasks at the same time (depends on your CPU
  cores).
\item
  Loops over all combinations of site\_ids and max\_irr\_values.
\item
  Collects results for all sites and irrigation levels into a list
  (all\_results).
\end{itemize}

\textbf{6. Convert results to a DataFrame and save}

results\_df = pd.DataFrame(all\_results)
results\_df.to\_csv(``./AquaCropReadme/merged\_results\_1.csv'',
index=False)

Converts the list of dictionaries into a DataFrame. Saves all the
optimized results in a CSV file for further analysis.

\begin{Shaded}
\begin{Highlighting}[]
\ImportTok{from}\NormalTok{ joblib }\ImportTok{import}\NormalTok{ Parallel, delayed}
\ImportTok{import}\NormalTok{ pandas }\ImportTok{as}\NormalTok{ pd}
\ImportTok{from}\NormalTok{ tqdm }\ImportTok{import}\NormalTok{ tqdm}

\CommentTok{\# {-}{-}{-}{-}{-}{-}{-}{-}{-}{-}{-}{-}{-}{-}{-}{-}{-}{-}{-}{-}{-}{-}{-}{-}{-}{-}{-}{-}{-}{-}{-}{-}{-}{-}{-}{-}{-}{-}{-}{-}{-}{-}{-}{-}}
\CommentTok{\# 1. Define irrigation levels to loop through}
\CommentTok{\# {-}{-}{-}{-}{-}{-}{-}{-}{-}{-}{-}{-}{-}{-}{-}{-}{-}{-}{-}{-}{-}{-}{-}{-}{-}{-}{-}{-}{-}{-}{-}{-}{-}{-}{-}{-}{-}{-}{-}{-}{-}{-}{-}{-}}
\NormalTok{max\_irr\_values }\OperatorTok{=}\NormalTok{ [}\DecValTok{10}\NormalTok{,}\DecValTok{20}\NormalTok{,}\DecValTok{30}\NormalTok{,}\DecValTok{40}\NormalTok{,}\DecValTok{50}\NormalTok{,}\DecValTok{60}\NormalTok{,}\DecValTok{70}\NormalTok{,}\DecValTok{80}\NormalTok{,}\DecValTok{90}\NormalTok{,}\DecValTok{100}\NormalTok{,}\DecValTok{110}\NormalTok{,}\DecValTok{120}\NormalTok{,}\DecValTok{130}\NormalTok{,}\DecValTok{140}\NormalTok{,}\DecValTok{150}\NormalTok{,}\DecValTok{160}\NormalTok{,}\DecValTok{170}\NormalTok{,}\DecValTok{180}\NormalTok{,}\DecValTok{190}\NormalTok{,}\DecValTok{200}\NormalTok{]}
\CommentTok{\# Modify as needed}

\CommentTok{\# {-}{-}{-}{-}{-}{-}{-}{-}{-}{-}{-}{-}{-}{-}{-}{-}{-}{-}{-}{-}{-}{-}{-}{-}{-}{-}{-}{-}{-}{-}{-}{-}{-}{-}{-}{-}{-}{-}{-}{-}{-}{-}{-}{-}}
\CommentTok{\# 2. Define site IDs}
\CommentTok{\# {-}{-}{-}{-}{-}{-}{-}{-}{-}{-}{-}{-}{-}{-}{-}{-}{-}{-}{-}{-}{-}{-}{-}{-}{-}{-}{-}{-}{-}{-}{-}{-}{-}{-}{-}{-}{-}{-}{-}{-}{-}{-}{-}{-}}
\NormalTok{site\_ids }\OperatorTok{=} \BuiltInTok{range}\NormalTok{(}\DecValTok{1}\NormalTok{, }\DecValTok{2}\NormalTok{)  }\CommentTok{\# Adjust as needed}

\CommentTok{\# {-}{-}{-}{-}{-}{-}{-}{-}{-}{-}{-}{-}{-}{-}{-}{-}{-}{-}{-}{-}{-}{-}{-}{-}{-}{-}{-}{-}{-}{-}{-}{-}{-}{-}{-}{-}{-}{-}{-}{-}{-}{-}{-}{-}}
\CommentTok{\# 3. Load all prepared weather files first}
\CommentTok{\# {-}{-}{-}{-}{-}{-}{-}{-}{-}{-}{-}{-}{-}{-}{-}{-}{-}{-}{-}{-}{-}{-}{-}{-}{-}{-}{-}{-}{-}{-}{-}{-}{-}{-}{-}{-}{-}{-}{-}{-}{-}{-}{-}{-}}
\NormalTok{wdfs }\OperatorTok{=}\NormalTok{ \{}
\NormalTok{    site\_id: prepare\_weather(}\SpecialStringTok{f"./ClimateData/site\_}\SpecialCharTok{\{}\NormalTok{site\_id}\SpecialCharTok{\}}\SpecialStringTok{.txt"}\NormalTok{)}
    \ControlFlowTok{for}\NormalTok{ site\_id }\KeywordTok{in}\NormalTok{ site\_ids}
\NormalTok{\}}

\CommentTok{\# {-}{-}{-}{-}{-}{-}{-}{-}{-}{-}{-}{-}{-}{-}{-}{-}{-}{-}{-}{-}{-}{-}{-}{-}{-}{-}{-}{-}{-}{-}{-}{-}{-}{-}{-}{-}{-}{-}{-}{-}{-}{-}{-}{-}}
\CommentTok{\# 4. Function that runs for each site + irrigation level}
\CommentTok{\# {-}{-}{-}{-}{-}{-}{-}{-}{-}{-}{-}{-}{-}{-}{-}{-}{-}{-}{-}{-}{-}{-}{-}{-}{-}{-}{-}{-}{-}{-}{-}{-}{-}{-}{-}{-}{-}{-}{-}{-}{-}{-}{-}{-}}
\KeywordTok{def}\NormalTok{ process\_site\_irr\_max(site\_id, max\_irr, wdf):}
\NormalTok{    smts }\OperatorTok{=}\NormalTok{ optimize(}\DecValTok{4}\NormalTok{, max\_irr, wdf)  }
\NormalTok{    yld, tirr, \_ }\OperatorTok{=}\NormalTok{ evaluate(smts, max\_irr, wdf, }\VariableTok{True}\NormalTok{)}

    \ControlFlowTok{return}\NormalTok{ \{}
        \StringTok{\textquotesingle{}Site\_ID\textquotesingle{}}\NormalTok{: site\_id,}
        \StringTok{\textquotesingle{}Max\_Irrigation\_mm\textquotesingle{}}\NormalTok{: max\_irr,}
        \StringTok{\textquotesingle{}Optimal\_SMTs\textquotesingle{}}\NormalTok{: smts.tolist(),}
        \StringTok{\textquotesingle{}Yield\_tonne\_per\_ha\textquotesingle{}}\NormalTok{: yld,}
        \StringTok{\textquotesingle{}Total\_Irrigation\_mm\textquotesingle{}}\NormalTok{: tirr}
\NormalTok{    \}}

\CommentTok{\# {-}{-}{-}{-}{-}{-}{-}{-}{-}{-}{-}{-}{-}{-}{-}{-}{-}{-}{-}{-}{-}{-}{-}{-}{-}{-}{-}{-}{-}{-}{-}{-}{-}{-}{-}{-}{-}{-}{-}{-}{-}{-}{-}{-}}
\CommentTok{\# 5. Run the parallel processing}
\CommentTok{\# {-}{-}{-}{-}{-}{-}{-}{-}{-}{-}{-}{-}{-}{-}{-}{-}{-}{-}{-}{-}{-}{-}{-}{-}{-}{-}{-}{-}{-}{-}{-}{-}{-}{-}{-}{-}{-}{-}{-}{-}{-}{-}{-}{-}}
\NormalTok{all\_results }\OperatorTok{=}\NormalTok{ Parallel(n\_jobs}\OperatorTok{=}\DecValTok{15}\NormalTok{)(}
\NormalTok{    delayed(process\_site\_irr\_max)(site\_id, max\_irr, wdfs[site\_id])}
    \ControlFlowTok{for}\NormalTok{ max\_irr }\KeywordTok{in}\NormalTok{ max\_irr\_values}
    \ControlFlowTok{for}\NormalTok{ site\_id }\KeywordTok{in}\NormalTok{ site\_ids}
\NormalTok{)}

\CommentTok{\# {-}{-}{-}{-}{-}{-}{-}{-}{-}{-}{-}{-}{-}{-}{-}{-}{-}{-}{-}{-}{-}{-}{-}{-}{-}{-}{-}{-}{-}{-}{-}{-}{-}{-}{-}{-}{-}{-}{-}{-}{-}{-}{-}{-}}
\CommentTok{\# 6. Convert results to dataframe and save}
\CommentTok{\# {-}{-}{-}{-}{-}{-}{-}{-}{-}{-}{-}{-}{-}{-}{-}{-}{-}{-}{-}{-}{-}{-}{-}{-}{-}{-}{-}{-}{-}{-}{-}{-}{-}{-}{-}{-}{-}{-}{-}{-}{-}{-}{-}{-}}
\NormalTok{results\_df }\OperatorTok{=}\NormalTok{ pd.DataFrame(all\_results)}

\NormalTok{results\_df.to\_csv(}
    \StringTok{"./AquaCropReadme/merged\_results\_1.csv"}\NormalTok{,}
\NormalTok{    index}\OperatorTok{=}\VariableTok{False}
\NormalTok{)}

\BuiltInTok{print}\NormalTok{(}\StringTok{"Finished! File saved."}\NormalTok{)}
\end{Highlighting}
\end{Shaded}

\begin{verbatim}
Finished! File saved.
\end{verbatim}

APPENDIX

Crop patemeters in AquaCrop

\begin{longtable}[]{@{}
  >{\raggedright\arraybackslash}p{(\columnwidth - 2\tabcolsep) * \real{0.5357}}
  >{\raggedright\arraybackslash}p{(\columnwidth - 2\tabcolsep) * \real{0.4643}}@{}}
\toprule\noalign{}
\begin{minipage}[b]{\linewidth}\raggedright
Variable Name
\end{minipage} & \begin{minipage}[b]{\linewidth}\raggedright
Description
\end{minipage} \\
\midrule\noalign{}
\endhead
\bottomrule\noalign{}
\endlastfoot
Name & Crop name, e.g., `maize' \\
CropType & Crop type (1 = Leafy vegetable, 2 = Root/tuber, 3 =
Fruit/grain) \\
PlantMethod & Planting method (0 = Transplanted, 1 = Sown) \\
CalendarType & Calendar type (1 = Calendar days, 2 = Growing degree
days) \\
SwitchGDD & Convert calendar to GDD mode if inputs are in calendar days
(0 = No; 1 = Yes) \\
PlantingDate & Planting date (mm/dd) \\
HarvestDate & Latest harvest date (mm/dd) \\
Emergence & Growing degree/calendar days from sowing to
emergence/transplant recovery \\
MaxRooting & Growing degree/calendar days from sowing to maximum
rooting \\
Senescence & Growing degree/calendar days from sowing to senescence \\
Maturity & Growing degree/calendar days from sowing to maturity \\
HIstart & Growing degree/calendar days from sowing to start of yield
formation \\
Flowering & Duration of flowering in growing degree/calendar days (-999
for non-fruit/grain crops) \\
YldForm & Duration of yield formation in growing degree/calendar days \\
GDDmethod & Growing degree day calculation method \\
Tbase & Base temperature (°C) below which growth does not progress \\
Tupp & Upper temperature (°C) above which crop development no longer
increases \\
PolHeatStress & Pollination affected by heat stress (0 = No, 1 = Yes) \\
Tmax\_up & Maximum air temperature (°C) above which pollination begins
to fail \\
Tmax\_lo & Maximum air temperature (°C) at which pollination completely
fails \\
PolColdStress & Pollination affected by cold stress (0 = No, 1 = Yes) \\
Tmin\_up & Minimum air temperature (°C) below which pollination begins
to fail \\
Tmin\_lo & Minimum air temperature (°C) at which pollination completely
fails \\
TrColdStress & Transpiration affected by cold temperature stress (0 =
No, 1 = Yes) \\
GDD\_up & Minimum GDD (°C/day) required for full crop transpiration
potential \\
GDD\_lo & GDD (°C/day) at which no crop transpiration occurs \\
Zmin & Minimum effective rooting depth (m) \\
Zmax & Maximum rooting depth (m) \\
fshape\_r & Shape factor describing root expansion \\
SxTopQ & Maximum root water extraction at top of root zone
(m³/m³/day) \\
SxBotQ & Maximum root water extraction at bottom of root zone
(m³/m³/day) \\
SeedSize & Soil surface area (cm²) covered by an individual seedling at
90\% emergence \\
PlantPop & Number of plants per hectare \\
CCx & Maximum canopy cover (fraction of soil cover) \\
CDC & Canopy decline coefficient (fraction per GDD/calendar day) \\
CGC & Canopy growth coefficient (fraction per GDD) \\
Kcb & Crop coefficient when canopy growth is complete but prior to
senescence \\
fage & Decline of crop coefficient due to ageing (\%/day) \\
WP & Water productivity normalized for ET₀ and CO₂ (g/m²) \\
WPy & Adjustment of water productivity in yield formation stage (\% of
WP) \\
fsink & Crop performance under elevated atmospheric CO₂ (\%/100) \\
HI0 & Reference harvest index \\
dHI\_pre & Possible increase of harvest index due to water stress before
flowering (\%) \\
a\_HI & Positive impact on HI of restricted vegetative growth during
yield formation \\
b\_HI & Negative impact on HI of stomatal closure during yield
formation \\
dHI0 & Maximum allowable increase of harvest index above reference
value \\
Determinant & Crop determinacy (0 = Indeterminant, 1 = Determinant) \\
exc & Excess of potential fruits \\
p\_up1 & Upper soil water depletion threshold for canopy expansion
stress \\
p\_up2 & Upper soil water depletion threshold for canopy stomatal
control stress \\
p\_up3 & Upper soil water depletion threshold for canopy senescence
stress \\
p\_up4 & Upper soil water depletion threshold for canopy pollination
stress \\
\end{longtable}

Irrigation Management strategies

Irrigation management parameters are selected by creating an
IrrigationManagement object. With this class we can specify a range of
different irrigation management strategies. The 6 different strategies
can be selected using the IrrMethod argument when creating the class.
These strategies are as follows:

\begin{itemize}
\tightlist
\item
  IrrMethod=0: Rainfed (no irrigation)
\item
  IrrMethod=1: Irrigation is triggered if soil water content drops below
  a specified threshold (or four thresholds representing four major crop
  growth stages (emergence, canopy growth, max canopy, senescence).
\item
  IrrMethod=2: Irrigation is triggered every N days
\item
  IrrMethod=3: Predefined irrigation schedule
\item
  IrrMethod=4: Net irrigation (maintain a soil-water level by topping up
  all compartments daily)
\item
  IrrMethod=5: Constant depth applied each day
\end{itemize}




\end{document}
